% ---------- TITLE PAGE ----------
\thispagestyle{empty}
\begin{tcolorbox}[enhanced, colback=hblue, colframe=hblue, boxrule=0pt, arc=4pt, left=16pt, right=16pt, top=22pt, bottom=22pt]
  \color{white}\sffamily
  {\large\bfseries Bewijzen in de Wiskunde \textperiodcentered\ 2026--2027}\\[10pt]
  {\fontsize{34}{38}\selectfont\bfseries Study Guide}\\[6pt]
  {\fontsize{22}{26}\selectfont Chapters 0--5: numbers, logic, sets, functions, induction, relations}\\[14pt]
  {\normalsize Theory you need for the \textbf{essential} and \textbf{recommended} exercises of the planning,\\
   with worked examples in the style the course expects.}
\end{tcolorbox}

\vspace{10pt}
\begin{tcolorbox}[enhanced, colback=white, colframe=hblue2, boxrule=0.8pt, arc=2pt, left=10pt, right=10pt, top=8pt, bottom=8pt]
\sffamily\small
\textbf{\textcolor{hblue}{Book:}} C.~Newstead, \emph{An Infinite Descent into Pure Mathematics}. Section numbers (\eg \S1.3, Theorem 5.2.35, exercise 0.E.22) refer to that book.\\[3pt]
\textbf{\textcolor{hblue}{Covers:}} Ch.\,0 (numbers, number bases, division, irrationals), \S1.1--1.3, \S2.1--2.2, \S3.1--3.2, \S4.2--4.3, \S5.1--5.2 and the chapter exercises 0.E--5.E named in the planning.\\[3pt]
\textbf{\textcolor{hblue}{Quizzes:}} Quiz 1 (23 Sep): numbers + propositional logic. Quiz 2 (7 Oct): sets + functions. Quiz 3 (21 Oct): relations + modular arithmetic. Exam 4 Nov (70\,\%).
\end{tcolorbox}

\vspace{8pt}
\hsub{How to use this guide}
Every chapter of this guide is one chapter of the book, and every lettered section (\lab{A}, \lab{B}, \dots) is one book section in the order of the planning. Each chapter runs up a ladder of three difficulty levels:

\begin{center}\sffamily\small
\begin{tabular}{@{}p{0.27\linewidth}p{0.68\linewidth}@{}}
\colorbox{hyellow}{\strut\textbf{Definition / key fact}} & Learn these word for word. Almost every exercise starts with ``unfold the definition''.\\[3pt]
\colorbox{hlight}{\strut\textbf{Theorem}} & Results you may quote in a proof (with their name or number).\\[3pt]
\colorbox{hgreen}{\strut\textbf{Proof strategy}} & The template to copy. The course grades the \emph{structure} of a proof as much as the maths.\\[3pt]
\colorbox{white}{\strut\textcolor{hblue}{\textbf{Example}} \textcolor{horange}{\textbf{Self Tutor}}} & \textbf{Level 1.} One proof shape at a time, written the way you should write it on a quiz. Cover the solution and try first.\\[3pt]
\colorbox{white}{\strut\textcolor{hpurple}{\textbf{Example}} \textcolor{hpurple}{\textbf{Book level}}} & \textbf{Level 2.} The difficulty of the hardest essential and recommended exercises: several shapes in one proof, abstract objects, and an \emph{invention step} that is named so you learn to spot it.\\[3pt]
\colorbox{white}{\strut\textcolor{hpurple}{\textbf{STRETCH EXERCISES}}} & \textbf{Level 3.} Unsolved problems with hints, at the top of the book's range. Full solutions in Appendix S. Do these before the exam, not before the quiz.\\[3pt]
\colorbox{hpurplebg}{\strut\textbf{EXERCISES}} & Which planning exercises the section prepares you for, and the skills they test.\\[3pt]
\colorbox{hredbg}{\strut\textbf{Watch out}} & The mistakes that cost the most marks.
\end{tabular}
\end{center}

Suggested order for each section: read the boxes, do the Level 1 examples with the solution covered, then the planning's essential exercises; then the Level 2 examples and the recommended exercises; then the stretch exercises when preparing for the exam.

\begin{key}[Conventions used in the book and in this guide]
\begin{itemize}
\item $\N=\{0,1,2,3,\dots\}$: \textbf{zero is a natural number} in this book. $\Z$, $\Q$, $\R$, $\C$ as usual.
\item A set is \term{inhabited} if it has at least one element (the book's word for ``non-empty'').
\item $X\subseteq Y$ means subset (equality allowed); $X\subsetneq Y$ means proper subset.
\item $f[U]$ is the image of $U\subseteq X$ and $f^{-1}[V]$ the preimage of $V\subseteq Y$ (square brackets; $f^{-1}[V]$ exists even when $f$ has no inverse).
\item $[a]_\sim$ or $[a]$ is the equivalence class of $a$; $[a]_n$ is the class of $a$ in $\Z/n\Z$.
\item Logical symbols: $\wedge$ and, $\vee$ or, $\Rightarrow$ implies, $\Leftrightarrow$ iff, $\neg$ not, $\forall$ for all, $\exists$ there exists, $\exists!$ there exists a unique.
\item $\square$ marks the end of a proof. \reason{grey text in braces} is the justification of a step.
\end{itemize}
\end{key}

\hsub{Planning at a glance (chapters 0--5)}
{\sffamily\small
\begin{tabularx}{\linewidth}{@{}l l X X@{}}
\toprule
\textbf{Date} & \textbf{Sections} & \textbf{Essential} & \textbf{Recommended}\\
\midrule
Mon 7 Sep & Ch.\,0 (to irrationals), A.1 & Ch0: 2, 21, 24, 40 \newline 0.E: 3, 4, 22, 27, 28 & Ch0: 14, 30, 38, 45, 51 \newline 0.E: 17, 18, 29, 30\\
Mon 14 Sep & \S1.1, \S1.3 (to ex.\,1.3.28) & 1.1: 21, 25, 32, 33, 38, 47, 56, 57 \newline 1.3: 5, 13, 14, 22, 27 & 1.1: 37, 43, 54, 55, 65, 66 \newline 1.E: 11\\
Wed 16 Sep & \S1.2, \S1.3 (from thm\,1.3.29) & 1.2: 6, 15, 18, 39 \newline 1.3: 44ace; 1.E: 18 & 1.2: 21, 22, 28, 37ab \newline 1.3: 32, 35; 1.E: 28, 29, 30\\
Mon 21 Sep & \S2.1, \S2.2 & 2.1: 10, 27, 31, 35 \newline 2.2: 5, 24, 26, 28, 35, 36, 37, 40, 51 \newline 2.E: 6, 33, 38 & 2.1: 17, 18, 23, 24, 32, 36 \newline 2.2: 4, 6, 9, 25, 41, 45, 53 \newline 2.E: 4, 5, 16, 28, 32, 41\\
Wed 23 Sep & \textbf{Quiz 1}; \S2.2, \S3.1 & 3.1: 11, 15abc, 36 \newline 3.E: 13, 17, 19, 34, 41 & 3.1: 8, 16, 19, 33, 40, 41 \newline 3.E: 2, 37, 38\\
Mon 28 Sep & \S3.2 & 3.2: 5, 8, 12, 14, 20, 40, 47 \newline 3.E: 24, 25, 48 & 3.2: 6, 16, 19, 25, 39ab, 45, 46 \newline 3.E: 35, 47\\
Wed 30 Sep & \S4.2 & 4.2: 5, 11 \newline 4.E: 10, 11, 16, 17, 18, 20, 25 & 4.2: 8, 17, 20 \newline 4.E: 12, 14, 15, 21, 24, 33\\
Mon 5 Oct & \S4.3 & 4.3: 8, 9, 14, 15 \newline 4.E: 13, 19, 26, 27, 29, 30, 31, 32 & 4.3: 10, 18 \newline 4.E: 28, 34, 35\\
Wed 7 Oct & \textbf{Quiz 2}; \S5.1, \S5.2 & 5.1: 6, 13, 22, 23, 28, 39, 40, 41 \newline 5.2: 24, 27, 33; 5.E: 18, 22, 23 & 5.1: 10, 11, 18, 24, 29, 33, 34 \newline 5.2: 25, 28, 30, 34, 36\\
\bottomrule
\end{tabularx}}

\vspace{6pt}
\hsub{Writing a proof (Appendix A.1 in one box)}
\begin{strategy}[The six habits of a readable proof]
\begin{enumerate}
\item \textbf{Say what you are proving and in which form.} ``We prove that for all $n\in\Z$, if $n$ is odd then $n^2$ is odd.'' Identify the outermost logical shape ($\forall$, $\Rightarrow$, $\wedge$, \dots): that shape decides the first line of the proof (Chapter 1).
\item \textbf{Introduce every variable before you use it.} ``Let $n\in\Z$.'' ``Fix $k\in\Z$ such that $n=2k+1$.'' A letter that appears from nowhere is an error.
\item \textbf{Unfold definitions.} ``$n$ is odd'' means ``$n=2k+1$ for some $k\in\Z$''. Most exercises are solved by unfolding the definitions on both sides and connecting them.
\item \textbf{Justify each step} with ``by definition of \dots'', ``by assumption'', ``by Theorem \dots'', or a one-line calculation. Never write ``clearly''.
\item \textbf{Write sentences.} $\Rightarrow$ and $\therefore$ are not abbreviations for ``so''. Chains of equations are fine; chains of unexplained symbols are not.
\item \textbf{Close the loop.} End with the statement you set out to prove and a $\square$. If the proof had cases, say that all cases are covered.
\end{enumerate}
\end{strategy}
